\section{Python IV reuse: rank and diagnostic names are insufficient contracts}
\label{exec22:pyiv-scope}

An instrument-based banking proposal must say which variation changes the
action, which exclusion restriction gives that variation a causal meaning,
and which sampling law justifies uncertainty. The pinned Python implementations
provide useful computational components, but their interfaces expose different
rank and diagnostic behavior. The selected \texttt{linearmodels} and
\texttt{pyfixest} revisions predate the literature cutoff; their equality to
particular distributed release binaries has not been established. The following
comparison concerns inspected source paths, with no native execution or claimed
coverage across the full Python IV and regression-discontinuity landscape
\citep{exec22_pyiv_kernels_linearmodels,exec22_pyiv_kernels_pyfixest}.

\subsection{The projected design and the causal exclusion restriction}
\label{exec22:pyiv-rank}

In the inspected \texttt{linearmodels} path, positive weights are normalized by
their mean, missing rows are removed, and square-root weights transform the
design, instrument matrix and outcome. The checks include instrument counts
and separate full column rank of $X$ and $Z$. For two-stage least squares the
solve uses the projected design, with matrix $X'P_ZX$, where
$P_Z=ZZ^+$ and $Z^+$ is the Moore--Penrose inverse. Separate ranks do not prove
that $X'P_ZX$ is invertible or well conditioned. They also say nothing about the
exclusion restriction~\citep{exec22_pyiv_kernels_linearmodels}.

An exact four-row witness makes both distinctions visible. Work with centred
vectors and no additional regressors, and define
\begin{equation}
 z=(-1,-1,1,1)',\qquad u=(-1,1,-1,1)',\qquad
 d=\varepsilon z+u,\qquad y=2d+z/100.
 \label{exec22:pyiv-witness}
\end{equation}
Here $z'u=0$ and $z'z=u'u=4$. Each single-column design $d$ and $z$ has full
column rank, including at $\varepsilon=0$. Nevertheless,
\begin{equation}
 d'P_zd=4\varepsilon^2,\qquad
 \widehat\beta_{\rm IV}=\frac{z'y}{z'd}
      =2+\frac{1}{100\varepsilon},\quad\varepsilon\ne0.
 \label{exec22:pyiv-ratio}
\end{equation}
At zero the projected design is singular. At $\varepsilon=1/100$ and $1/1000$
the fitted ratios are $3$ and $12$, while the stipulated action coefficient is
$2$. The term $z/100$ represents a direct instrument effect and violates
exclusion. For a branch encouragement instrument it could represent a small
direct change in service attention outside the action being evaluated.
Small direct effects become consequential when the first stage is weak.
This is an exact sensitivity witness, not a sampling experiment, a package
failure or an empirical assertion about a bank~\citep{exec22_pyiv_reference}.

The inspected \texttt{pyfixest} path removes collinear within-transformed
columns, retaining their names and a mask, before delegating weighted 2SLS.
Its selected fit kernel constructs the first-stage projection with $Z'WZ$,
fits the projected design, retains residuals in response units and includes
weights in the score contributions. The collinearity wrapper passes the
unweighted within Gram matrix to a separate core routine.
Section~\ref{sec:exec25-python-iv-rd} follows that routine and the solver
backend, finding an order- and tolerance-dependent column test and solver
branches that do not retain every returned diagnostic. These findings make
it necessary to retain the actual surviving variables, weight roles and
tolerances at the interface; native numerical qualification remains open. It is not evidence that a column-removal warning is a successful
causal repair~\citep{exec22_pyiv_kernels_pyfixest}.

\subsection{Preserve the actual covariance model}
\label{exec22:pyiv-covariance}

The selected \texttt{linearmodels} fit defaults to heteroskedastic-robust
covariance with the finite-sample adjustment disabled. Its robust score meat
uses row contributions and does not by itself account for serial dependence.
The kernel path has a separate bandwidth choice, including an automatically
selected option. The clustering class documents one or two grouping columns
and contains a two-way inclusion--exclusion branch, even though a nearby class
name and fit description say one-way. Its two-way score covariance takes the
form $S_1+S_2-S_{12}$ with component adjustments when requested. Actual grouping
columns and the intersection construction must therefore accompany the result;
a label alone does not establish the implemented dependence model
\citep{exec22_pyiv_kernels_linearmodels}.

In the absence of an explicit grouping array, the inspected class assigns
each row its own cluster. That is a computational default, not a choice of
bank-level or calendar-level dependence. An adapter should require the intended
sampling units and ordering, report the realized HAC bandwidth or cluster counts,
and retain covariance warnings or singularity. Additional validation is needed
for small numbers of groups, joint calendar dependence, group-code conversion
and the covariance's statistical assumptions. No uniform robustness to arbitrary
banking dependence is implied.

\subsection{A statistic and its p-value can use different displayed scalings}
\label{exec22:pyiv-diagnostics}

For a restriction $R\beta=q$, distinguish the Wald quadratic $W$ from its
F-scaled display $W/k$, where $k$ is the number of restrictions. In the selected
\texttt{pyfixest} parent method, an F reference is used only for the joint zero
restriction on all coefficients. A subset restriction changes the reference
to chi-square and emits a warning. The first-stage helper selects excluded
instruments, suppresses warnings around this call, then retains the F-scaled
value and the returned p-value. When excluded instruments are a proper subset
of the first-stage coefficients, these fields can therefore combine $W/k$
with a chi-square p-value. This is a source-level description of the selected
branch, not evidence of an executed diagnostic error
\citep{exec22_pyiv_kernels_pyfixest}.

For example, $W=8$ and $k=2$ give displayed $W/k=4$. A chi-square-two reference
gives $p=e^{-4}\simeq0.0183$, whereas an $F_{2,20}$ reference gives
$p=(1+4/10)^{-10}\simeq0.0346$. The same displayed number cannot identify
which p-value is intended. The selected low-level Wald utility retains both
scalings and the distribution, but the summarized first-stage object drops
some of that information. A wrapper should recover the parent result and
retain the null, restriction matrix, covariance, actual distribution, degrees
of freedom, both scalings and the p-value. A covariance pseudoinverse in the
quadratic calculation does not itself justify inference from singular or
invalid uncertainty~\citep{exec22_pyiv_kernels_pyfixest,exec22_pyiv_reference}.

The earlier \texttt{linearmodels} diagnostic inspection likewise distinguishes
its first-stage F or chi-square reference and its \texttt{anderson\_rubin}
overidentification property~\citep{exec22_pyiv_diagnostics_linearmodels}.
The latter's method name is not a contract for
testing every proposed coefficient value. A weak-IV confidence set requires
the appropriate coefficient-null test and its inversion, including empty,
disconnected or unbounded results when justified. The proposed CLI and MCP
surfaces should expose named statistical operations with those targets, rather
than map similarly named fields to a common confidence-interval slot.

For prototype selection, preserve these source distinctions and compare
matched models before choosing either engine. Qualify row and column conversion,
weight semantics, projection conditioning, covariance and diagnostic metadata
before comparing speed. Generated designs should separately vary instrument
strength, exclusion violations, dependence and retained samples, with failures
counted in the denominator. Section~\ref{sec:exec25-python-iv-rd} supplies the
separate Python RD and CausalPy IV source comparison, including model and prior
contracts, cutoff conventions and confidence-set geometry. Those findings
extend this kernel analysis without granting general landscape completion or
native admission. The next investment should buy a
specific testable capability, with its unresolved assumptions visible to the
decision maker.
